% part: intuitionistic-logic
% chapter: propositions-as-types
% section: normalization

\documentclass[../../../include/open-logic-section]{subfiles}

\begin{document}

\olfileid{int}{pty}{nor}

\olsection{نارمل روپ وچ لیاؤنا}

ایس حصے وچ اسیں ثابت کردے آں کہ تقلیل دی کسے ترتیب نال ہر اخذ نوں نارمل اخذ وچ گھٹایا جا سکدا اے۔ ایس نوں \emph{نارمل روپ وچ لیاؤن دی خاصیت} آکھدے آں۔ اسیں قضیاں تے نوعاں دا میل ورتاں گے۔ اسیں وکھاؤنے آں کہ ہر ثبوت ٹرم نوں نارمل روپ وچ گھٹایا جا سکدا اے۔ ایہدے توں قدرتی اخذ دے !!{derivation}s لئی نارمل روپ وچ لیاؤن دا نتیجہ نکلدا اے۔

پہلاں اسیں کجھ فنکشناں دی تعریف دیندے آں جیہڑے ٹرماں دی گنجل ماپدے نیں۔ \emph{لمبائی}~$\len{!A}$، !!a{formula}s دی، ایویں متعین اے:
\begin{align*}
  \len{p} &= 0\\
  \len{!A \land !B} &= \len{!A} + \len{!B} + 1\\
  \len{!A \lor !B} &= \len{!A} + \len{!B} + 1 \\
  \len{!A \lif !B} &= \len{!A} + \len{!B} + 1.
\end{align*}
ریڈیکس~$M$ دی گنجل اوہدے \emph{قطع درجے}~$\cutrank{M}$ نال ماپی جاندی اے:
\begin{align*}
  \cutrank{(\lambd[\typeof{x}{!A}][\typeof{N}{!B}])Q} &=
  \len{!A} + \len{!B} + 1 \\
  \cutrank{\ande{i}{\andi{\typeof{M}{!A}}{\typeof{N}{!B}}}} &=
  \len{!A} + \len{!B} + 1 \\
  \cutrank{\ore{\ori{i}{!A_{i}}{\typeof{M}{!A_i}}}
    {\typeof{x_1}{!A_1}}{\typeof{N_1}{!C}}
    {\typeof{x_2}{!A_2}}{\typeof{N_2}{!C}}} &=
  \len{!A} + \len{!B} + 1
\end{align*}
ثبوت ٹرم دی گنجل اوہدے وچلے سبھ توں گنجل دار ریڈیکس نال ماپی جاندی اے، تے جے ٹرم نارمل اے، تاں $0$ اے:
\[
  \maxrank{M} = \max \{\cutrank{N} | N \text{ ذیلی ٹرم اے } M \text{ دا تے ریڈیکس اے}\}
\]

\begin{lem}\ollabel{lem:subst}
  جے $\Subst{M}{\typeof{N}{!A}}{\typeof{x}{!A}}$ ریڈیکس اے تے $M
  \not\ident x$، تاں ہِٹھلیاں صورتاں وچوں اک اے:
  \begin{enumerate}
  \item $M$ آپ ریڈیکس اے، یا
  \item $M$ دی شکل $\ande{i}{x}$ اے، تے $N$ دی شکل $\andi{P_1}{P_2}$ اے۔
  \item $M$ دی شکل $\ore{i}{x_1}{P_1}{x_2}{P_2}$ اے، تے $N$ دی شکل $\ori{i}{}{Q}$ اے۔
  \item $M$ دی شکل $x Q$ اے، تے $N$ دی شکل $\lambd[x][P]$ اے۔
  \end{enumerate}
  پہلی صورت وچ $\cutrank{\Subst{M}{N}{x}} = \cutrank{M}$، تے ہورناں صورتاں وچ $\cutrank{\Subst{M}{N}{x}} = \len{!A})$۔
\end{lem}
\begin{proof}
  ثبوت~$M$ اُتے استقرا نال اے۔
  \begin{enumerate}
  \item جے $M$ اکّو متغیر $y$ اے تے $y \not\ident x$، تاں $\Subst{y}{N}{x}$، $y$ اے، ایس لئی ریڈیکس نہیں۔
  \item جے $M$ دی شکل $\andi{N_1}{N_2}$، یا $\lambd[x][N]$، یا $\ori{i}{!A}{N}$ اے، تاں $\Subst{M}{\typeof{N}{!A}}{\typeof{x}{!A}}$ وی اوسے شکل دا اے، ایس لئی ریڈیکس نہیں۔
  \item جے $M$ دی شکل $\ande{i}{P}$ اے، تاں دو صورتاں ویکھدے آں۔
    \begin{enumerate}
      \item جے $P$ دی شکل $\andi{P_1}{P_2}$ اے، تاں $M \ident
        \ande{i}{\andi{P_1}{P_2}}$ ریڈیکس اے، تے صاف اے کہ
        \[
        \Subst{M}{N}{x} \ident
        \ande{i}{\andi{\Subst{P_1}{N}{x}}{\Subst{P_2}{N}{x}}}
        \]
        وی ریڈیکس اے۔ قطع درجے برابر نیں۔
      \item جے $P$ اکّو متغیر اے، تاں تھانواں بدلّن نال ریڈیکس بنن لئی اوہ $x$ ہونا چاہیدا اے، تے $N$ دی شکل $\andi{P_1}{P_2}$ ہونی چاہیدی اے۔ ہُن ایہ ویکھو:
        \[
        \Subst{M}{N}{x} \ident \Subst{\ande{i}{x}}{\andi{P_1}{P_2}}{x},
        \]
        جیہڑا $\ande{i}{\andi{P_1}{P_2}}$ اے۔ اوہدا قطع درجہ $\cutrank{x}$ دے برابر اے، جیہڑا $\len{!A}$ اے۔
    \end{enumerate}
  \end{enumerate}
  $\ore{N}{x_1}{N_1}{x_2}{N_2}$ تے $P Q$ دیاں صورتاں ایہو جیہیاں نیں۔
\end{proof}

\begin{lem}
  جے $M$ سکڑ کے $M'$ بن جاندا اے، تے $\cutrank{M} > \cutrank{N}$، ہر ریڈیکس ذیلی ٹرم $N$ لئی جیہڑا $M$ توں وکھرا اے، تاں $\cutrank{M} >
  \maxrank{M'}$۔
\end{lem}

\begin{proof}
  وکھ وکھ صورتاں نال ثبوت۔
  \begin{enumerate}
  \item جے $M$ دی شکل $\ande{i}{\andi{M_1}{M_2}}$ اے، تاں $M'$، $M_i$ اے۔ کیونکہ $M_i$ دا ہر ذیلی ٹرم $M$ دا وی اوہدے اپنے توں وکھرا ذیلی ٹرم اے، ایس لئی دعویٰ ٹھیک اے۔
  \item جے $M$ دی شکل $(\lambd[\typeof{x}{!A}][N])\typeof{Q}{!A}$ اے، تاں $M'$، $\Subst{N}{\typeof{Q}{!A}}{\typeof{x}{!A}}$ اے۔ $M'$ وچ کسے ریڈیکس نوں ویکھو۔ یا $N$ وچ اوہدے نال میل کھاندا اوہو قطع درجے والا ریڈیکس اے، جیہڑا مفروضے دے حساب نال $\cutrank{M}$ توں گھٹ اے، یا قطع درجہ $\len{!A}$ دے برابر اے، جیہڑا تعریف دے حساب نال $\cutrank{(\lambd[x^{!A}][N])Q}$ توں گھٹ اے۔
  \item جے $M$ دی شکل
    \[
    \ore{\ori{i}{}{\typeof{N}{!A_i}}}
        {\typeof{x_1}{!A_1}}{\typeof{N_1}{!C}}
        {\typeof{x_2}{!A_2}}{\typeof{N_2}{!C}},
    \]
    اے، تاں $M' \ident \Subst{N_i}{N}{\typeof{x_i}{!A_i}}$۔~$M'$ وچ کسے ریڈیکس نوں ویکھو۔ یا $N_i$ وچ اوہدے نال میل کھاندا اوہو قطع درجے والا ریڈیکس اے، جیہڑا مفروضے دے حساب نال $\cutrank{M}$ توں گھٹ اے، یا قطع درجہ $\len{!A_i}$ دے برابر اے، جیہڑا تعریف دے حساب نال $\cutrank{\ore{\ori{i}{}{\typeof{N}{!A_i}}}
      {\typeof{x_1}{!A_1}}{\typeof{N_1}{!C}}
      {\typeof{x_2}{!A_2}}{\typeof{N_2}{!C}}}$ توں گھٹ اے۔
  \end{enumerate}
\end{proof}

\begin{thm}
  سارے ثبوت ٹرم تقلیل نال نارمل روپ وچ آ جاندے نیں۔ سارے !!{derivation}s تقلیل نال نارمل !!{derivation}s بن جاندے نیں۔
\end{thm}

\begin{proof}
  دوجی گل پہلی توں نکلدی اے۔ اسیں پہلی گل $m=\maxrank{M}$ اُتے پورے استقرا نال ثابت کردے آں، جتھے $M$ ثبوت ٹرم اے۔
  \begin{enumerate}

  \item جے $m=0$، تاں $M$ پہلاں ای نارمل اے۔
  \item نئیں تاں اسیں~$n$ اُتے استقرا کردے آں۔ ایہ $M$ وچ اوہناں ریڈیکساں دی گنتی اے جیہناں دا قطع درجہ~$m$ دے برابر اے۔
    \begin{enumerate}
    \item جے $n=1$، تاں کوئی ریڈیکس $N$ چُنو ایویں کہ $m = \cutrank{N} >
      \cutrank{P}$، ہر اوہدے اپنے توں وکھرے ذیلی ٹرم $P$ لئی جیہڑا آپ وی ریڈیکس اے۔ ایہو جیہا ریڈیکس ضرور موجود اے، کیونکہ ہر ٹرم دے ذیلی ٹرم محدود گنتی وچ نیں۔

      منّ لو کہ $N'$، $N$ دی تقلیل نال ملن والا ٹرم اے۔ ہُن لیمے نال $\maxrank{N'} < \maxrank{N}$، ایس لئی اسیں ویکھدے آں کہ $n$، یعنی $\cutrank=m$ والے ریڈیکساں دی گنتی، گھٹ گئی اے۔ سو $m$ وی گھٹ گیا اے ($1$ یا ودھ نال)، تے اسیں~$m$ لئی استقرا دا مفروضہ لا سکدے آں۔
    \item استقرائی قدم لئی منّ لو $n>1$۔ طریقہ رلدا ملدا اے، فرق ایہ اے کہ $n$ گھٹ کے وی دھناں عدد رہندا اے، سو $m$ نہیں بدلدا۔ اسیں~$n$ لئی استقرا دا مفروضہ لا دیندے آں۔
    \end{enumerate}
\end{enumerate}
\end{proof}

نوع بندی والے $\lambd$-ٹرم نارمل روپ وچ آ سکدے نیں، یعنی ہر ایہو جیہے ٹرم دا \emph{اک نارمل روپ اے}۔ جے $\beta$-تقلیل نال $\lambd$-ٹرماں نوں نارمل روپ وچ لیاؤن نوں حسابی عمل چلاؤنا تے نارمل روپ نوں اوس عمل دی قدر سمجھیے، تاں ایہدا مطلب اے کہ نوع بندی والے $\lambd$-کیلکولس وچ ہر حسابی عمل آخر رُک جاندا اے تے اک قدر دیندا اے۔ نالے نوع دی بقا پکی کردی اے کہ قدر دی نوع ٹھیک اے۔ مثال لئی جے $N$ دی نوع $!A \lif !B$ اے (یعنی اوہ ایہو جیہا فنکشن اے جیہڑا~$!A$ نوع دے آرگومنٹاں اُتے کم کرے تے~$!B$ نوع دیاں قدراں دیوے)، تے اسیں اوہنوں~$M$ ٹرم اُتے لائیے جیہدی نوع~$!A$ اے (ان پُٹ)، تاں $(NM)$ گھٹ کے ٹرم $N'$ بن جاندا اے (آؤٹ پُٹ)، تے پکا اے کہ اوس آؤٹ پُٹ دی نوع~$!B$ اے۔

اصل وچ نوع بندی والے $\lambd$-کیلکولس دے ٹرم صرف اک طریقے نال ذیلی ٹرماں اُتے $\beta$-تقلیل لا کے نارمل روپ وچ نہیں آؤندے۔ ذیلی ٹرماں اُتے $\beta$-تقلیل لاؤن دا \emph{ہر} طریقہ آخر نارمل روپ اُتے مُکدا اے۔ ایس خاصیت نوں \emph{مضبوط نارملائزیشن} آکھدے آں۔ ایہدا مطلب اے کہ حسابی عمل نوں آخر روک کے قدر حاصل کرن دی ضمانت والا طریقہ موجود اے، تے حسابی عمل چلاؤن دا \emph{ہر} طریقہ آخر قدر دیندا اے۔

سانوں کیویں پتا اے کہ حسابی عمل چلاؤن دے وکھ وکھ طریقے \emph{اوہی} قدر دیندے نیں؟ ہُن تک سانوں بس ایہ پتا اے (نوع دی بقا توں) کہ حاصل ہوئیاں ساریاں قدراں اکّو نوع دیاں نیں۔ نوع بندی والے $\lambd$-کیلکولس دی ہور اہم خاصیت اے۔ ایہ \emph{سنگم والی} یا \emph{چرچ–راسر} خاصیت رکھدا اے۔ جے $N \red N_1$ تے $N \red
N_2$، تاں کوئی ٹرم~$N'$ اے ایویں کہ $N_1 \red N'$ تے $N_2 \red
N'$۔ چیتے رکھو کہ $N \red N'$ دا مطلب ایہ اے کہ $N'$ صفر یا ودھ $\redone$ تقلیل قدماں نال ملدا اے۔ جے $N_1$ نارمل روپ وچ اے، تاں اکّو ٹرم $N'$ جیہدے لئی $N_1 \red N'$ اے، $N_1$ آپ اے (صفر $\redone$ تقلیل قدماں نال)۔ ایس توں نکلدا اے کہ جے $N_1$ تے $N_2$ دونے نارمل روپ وچ نیں، تاں $N_1 = N' = N_2$۔ دوجے لفظاں وچ، کیونکہ $\red$ مضبوط نارملائزیشن والی اے، ٹرم دی تقلیل دا ہر طریقہ آخر نارمل روپ دیندا اے۔ کیونکہ $\red$ چرچ–راسر اے، کوئی وی دو نارمل روپ برابر نیں۔ ایس لئی نوع بندی والے $\lambd$-کیلکولس وچ حسابی عمل ہمیشہ مُک جاندا اے، تے اوہدی قدر (نارمل روپ) اوہی رہندی اے بھاویں حسابی عمل کِسے وی طرح چلے۔

ایہ ثابت کرنا عام طور اُتے اوکھا اے کہ کسے تعلق دی سنگم والی خاصیت اے۔ پر اسیں ہِٹھلی کمزور شرط سوکھے ای ثابت کر سکدے آں۔

\begin{prop}[چرچ–راسر دی کمزور خاصیت]
جے $N \redone N_1$ تے $N \redone N_2$، تاں اک ٹرم $N'$ اے ایویں کہ $N_1 \red N'$ تے $N_2 \red N'$۔
\end{prop}

\begin{proof}
  دونے ٹرم $N_1$ تے $N_2$، $N$ توں ملدے نیں، بالترتیب ریڈیکس $M_1$ یا $M_2$ دی تھاں اوہدی تقلیل نال ملن والا ٹرم رکھ کے۔ ریڈیکس $M_1$ تے $M_2$، $N$ دے ذیلی ٹرم نیں، سو اوہ اک دوجے اُتے نہیں چڑھ سکدے۔ ایہدا مطلب اے کہ یا اک دوجے دا ذیلی ٹرم اے، یا اوہ $N$ وچ دو وکھریاں، سانجھ توں بناں تھاواں اُتے نیں۔ دوجی صورت ویکھو۔ $N$ دی شکل $N(M_1,M_2)$ اے۔ $M_1$ دی تھاں اوہدی تقلیل نال ملن والا $M_1'$ رکھن نال $N_1$ ملدا اے، جیہڑا $N(M_1',M_2)$ اے۔ تے $M_2$ دی تھاں اوہدی تقلیل نال ملن والا $M_2'$ رکھن نال $N_2$ ملدا اے، جیہڑا $N(M_1,M_2')$ اے۔ دونے گھٹ کے $N(M_1',M_2')$ بن جاندے نیں۔ ہور صورت لئی منّ لو $M_2$،~$M_1$ دا ذیلی ٹرم اے (دوجی صورت ایہدے نال سم روپی اے)، تے $M_1$ کِس طرح دا ریڈیکس اے، اوس دے حساب نال صورتاں وکھ کرو۔ مثال لئی جے $M_1 \ident
  \proj{1}{\pair{O_1}{O_2}}$، تاں اوہ گھٹ کے $O_1$ بن جاندا اے۔ جے $M_2$، $O_1$ دا ذیلی ٹرم اے، تاں $M_2$ دی تبدیلی نال $O_1'$ ملدا اے۔ سو پہلاں $M_1$ تے فیر $M_2$ بدلّن نال $O_1'$ ملدا اے۔ پر پہلاں $M_2$ بدلّن نال $\proj{1}{\pair{O_1'}{O_2}}$ ملدا اے تے ایہ وی گھٹ کے~$O_2$ بن جاندا اے۔
\end{proof}

\begin{lem}[نیومین دا لیما]
جے $\red$ مضبوط نارملائزیشن والی اے تے چرچ–راسر دی کمزور خاصیت رکھدی اے، تاں اوہ سنگم والی اے۔
\end{lem}

\begin{proof}
  کیونکہ $\red$ مضبوط نارملائزیشن والی اے، ہر تقلیل لڑی اک نارمل روپ دیندی اے۔ نارمل روپ اوہدوں تے صرف اوہدوں اکّو ہوندا اے جدوں $\red$ سنگم والی اے۔ سو منّ لو کہ کسے ٹرم $N$ دے دو وکھرے نارمل روپ $N'$ تے $N''$ نیں، جیہڑے تقلیل لڑیاں توں ملے نیں:
  \begin{align*}
  N & = N_1' \redone N_2' \redone \dots \redone N_k' = N' \text{ تے}\\
  N & = N_1'' \redone N_2'' \redone \dots \redone N_l'' = N''
  \end{align*}
  (تقلیل لڑی دی لمبائی $>0$ ہونی چاہیدی اے، نئیں تاں $N$ پہلاں ای نارمل روپ وچ ہووے گا۔)

  جے $N_1' = N_1''$، تاں ایہدے دو وکھرے نارمل روپ نیں ($N'$ تے $N''$)۔ جے $N_1' \neq N_1''$، تاں چرچ–راسر دی کمزور خاصیت نال کوئی $N'''$ اے ایویں کہ $N_1' \red N'''$ تے $N_1'' \red N'''$۔ کیونکہ $N' \neq N''$، تاں یا $N''' \neq N'$ یا $N''' \neq N''$۔
  ایس لئی $N_1'$ یا $N_1''$ وچوں کسے اک دے دو وکھرے نارمل روپ نیں۔ اسیں وکھایا اے کہ جے $N$ دے دو وکھرے نارمل روپ نیں، تاں اک $N_1$ اے ایویں کہ $N \redone N_1$ تے $N_1$ دے دو وکھرے نارمل روپ نیں۔ صاف اے کہ اسیں ایویں اگے ودھ سکدے آں تے $N \redone N_1
  \redone N_2 \redone \dots$ بنا سکدے آں، پر ایہ ممکن نہیں، کیونکہ $\redone$ مضبوط نارملائزیشن والی اے۔
\end{proof}

\end{document}
